% !TeX program = xelatex
\documentclass[12pt,spanish,letterpaper,oneside]{article}
\PassOptionsToPackage{es-noshorthands}{babel}
\def\documenttitle{Práctica reflexiva sobre dilemas éticos profesionales y cultura de la legalidad en el ejercicio agronómico}
\def\documentsubtitle{Evidencia 2 --- Informe de práctica reflexiva y deontología aplicada (Fase 2)}
\def\documentsubject{Ética y Cultura de la Legalidad --- Ingeniero Agrónomo}
\def\documentauthor{Carlos Mauricio Carvajal Coronado}
\def\coursename{Ética y Cultura de la Legalidad}
\def\coursecode{REF\#261013755115}
\def\universityname{Universidad Autónoma de Nuevo León}
\def\universityfaculty{Facultad de Agronomía}
\def\universitydepartment{Ingeniero Agrónomo}
\def\universitydepartmentimage{departamentos/logos-institucionales.pdf}
\def\universitydepartmentimagecfg{height=2.1cm}
\def\universitylocation{General Escobedo / Marín, Nuevo León, México}
\def\authortable{%
	\begin{Form}
	\begin{tabular}{@{}ll@{}}
		\textbf{Estudiante:} & Carlos Mauricio Carvajal Coronado\\[0.35cm]
		\textbf{Matrícula:} & 1977875\\[0.35cm]
		\textbf{Docente:} & Cátedra de Formación General Universitaria\\[0.35cm]
		\textbf{Unidad de Aprendizaje:} & Ética y Cultura de la Legalidad (Evidencia 2 \textbar\ 8 Pts)\\[0.5cm]
		\multicolumn{2}{l}{Fecha de realización: 14 de octubre de 2026}\\
		\multicolumn{2}{l}{\universitylocation}
	\end{tabular}
	\end{Form}
}

\usepackage{fontspec}
\setmainfont{Arial}
\setsansfont{Arial}
\AddToHook{file/config.tex/after}[aulatex-etica]{%
  \def\documentfontsize{12}
  \def\documentinterline{1}
  \def\documentparindent{0}
  \def\documentparskip{5}
  \def\fontdocument{custom}
  \def\fonttypewriter{custom}
  \def\stylecitereferences{custom}
  \def\portraitstyle{style1}
  \def\hfstyle{style1}
  \def\hfwidthwrap{true}
  \def\hfwidthtitle{0.44}
  \def\hfwidthcourse{0.52}
  \def\predocpageromannumber{false}
  \def\pagemarginleft{2.5}
  \def\pagemarginright{2.5}
  \def\pagemargintop{2.5}
  \def\pagemarginbottom{2.5}
  \def\firstpagemargintop{3.8}
  \def\sectionfontsize{\large}
  \def\sectionspacingtop{8}
  \def\sectionspacingbottom{4}
  \def\ssectionfontsize{\normalsize}
  \def\ssectionspacingtop{6}
  \def\ssectionspacingbottom{3}
  \def\defaultimagefolder{UANL/img/}
  \def\usepdfmetadata{false}
}
% Template:     Informe LaTeX
% Documento:    Núcleo del template
% Versión:      DEV
% Codificación: UTF-8
%
% Autor: Pablo Pizarro R.
%        pablo@ppizarror.com
%
% Manual template: [https://latex.ppizarror.com/informe]
% Licencia MIT:    [https://opensource.org/licenses/MIT]

% -----------------------------------------------------------------------------
% CONFIGURACIONES
% -----------------------------------------------------------------------------
% Definiciones previas
\input{src/defs}

% Ajustes usuario
\input{src/config}

% -----------------------------------------------------------------------------
% IMPORTACIÓN DE LIBRERÍAS
% -----------------------------------------------------------------------------
\input{src/env/imports}

% -----------------------------------------------------------------------------
% IMPORTACIÓN DE FUNCIONES Y ENTORNOS
% -----------------------------------------------------------------------------
\input{src/cmd/core}
\input{src/cmd/math}
\input{src/cmd/equation}
\input{src/cmd/image}
\input{src/cmd/title}
\input{src/cmd/other}
\input{src/cmd/column}
\input{src/env/environments}

% -----------------------------------------------------------------------------
% IMPORTACIÓN DE ESTILOS
% -----------------------------------------------------------------------------
\input{src/style/code}
\input{src/style/other}

% -----------------------------------------------------------------------------
% CONFIGURACIÓN INICIAL DEL DOCUMENTO
% -----------------------------------------------------------------------------
\input{src/cfg/init}

% -----------------------------------------------------------------------------
% PORTADA
% -----------------------------------------------------------------------------
\input{src/page/portrait}
\input{src/page/portrait_config}

% -----------------------------------------------------------------------------
% CONFIGURACIÓN DE PÁGINA Y ENCABEZADOS
% -----------------------------------------------------------------------------
\input{src/cfg/page}

% -----------------------------------------------------------------------------
% TABLA DE CONTENIDOS - ÍNDICE
% -----------------------------------------------------------------------------
\input{src/page/index}

% -----------------------------------------------------------------------------
% CONFIGURACIONES FINALES
% -----------------------------------------------------------------------------
\input{src/cfg/final}

\RemoveFromHook{file/config.tex/after}[aulatex-etica]


\usepackage{tabularx}
\usepackage{booktabs}
\usepackage{amsmath,amssymb}
\usepackage[backend=biber,style=apa,sorting=nyt]{biblatex}
\definecolor{uanlblue}{RGB}{0,55,108}
\definecolor{tableblue}{RGB}{233,240,247}
\def\portraittitlecolor{uanlblue}
\def\sectioncolor{uanlblue}
\def\ssectioncolor{uanlblue}
\def\documenttitlehf{Ética y Legalidad: Evidencia 2 --- Práctica Reflexiva Deontológica}
\hypersetup{colorlinks=true,linkcolor=uanlblue,urlcolor=uanlblue,citecolor=uanlblue,pdfauthor={Carlos Mauricio Carvajal Coronado},pdftitle={Práctica reflexiva sobre dilemas éticos en agronomía},pdfsubject={Evidencia 2 - Ética y Legalidad UANL}}
\addbibresource{UANL/ingeniero-agronomo/etica-y-cultura-de-la-legalidad/etica-y-cultura-de-la-legalidad.bib}
\newcolumntype{Y}{>{\raggedright\arraybackslash}X}
\setlength{\emergencystretch}{2.5em}

\begin{document}
\hypersetup{pageanchor=false}
\begingroup
\fontsize{14}{17}\selectfont
\renewcommand{\Huge}{\fontsize{14}{17}\selectfont\bfseries}
\renewcommand{\large}{\fontsize{14}{17}\selectfont}
\templatePortrait
\clearpage
\endgroup
\hypersetup{pageanchor=true}
\templatePagecfg
\singlespacing
\setlength{\tabcolsep}{8pt}
\renewcommand{\arraystretch}{1.18}

\section*{Resumen}
En el presente trabajo se formaliza la Evidencia 2 de la Unidad de Aprendizaje de Ética y Cultura de la Legalidad mediante un Informe de Práctica Reflexiva orientado a dilucidar los dilemas éticos y normativos a los que se enfrenta el Ingeniero Agrónomo en el mercado profesional del noreste de México. Se analiza de forma crítica la disyuntiva entre presiones de rentabilidad comercial y el cumplimiento estricto del marco fitosanitario, evaluando casos de prescripción de moléculas plaguicidas no autorizadas o con periodos de carencia violentados, y la comercialización de germoplasma no certificado. Con apoyo en los postulados de Miguel Giusti sobre deliberación moral, la teoría de la responsabilidad de Hans Jonas y el Código de Ética de la UANL, se delimitan los criterios de discernimiento ético, la prevalencia del Estado de derecho y los límites morales de la lealtad laboral frente al bienestar público.

\vspace{0.4cm}
\noindent\textbf{Palabras clave:} Práctica reflexiva, dilemas éticos, deontología agronómica, sanidad vegetal, agroquímicos, Código de Ética UANL.

\section{Introducción}
El ejercicio de la agronomía conlleva una trascendental función social: la garantía de la inocuidad alimentaria, la preservación ecológica y la justicia distributiva rural \parencite{cortina2013}. No obstante, en la práctica profesional cotidiana emergen tensiones entre el mandato ético-deontológico y las exigencias de maximización de ganancias de productores o casas agroquímicas \parencite{giusti2018}. El análisis reflexivo sistemático de estos dilemas constituye un ejercicio formativo esencial para consolidar una cultura de la legalidad inquebrantable en la Facultad de Agronomía de la UANL.

\section{Identificación y Análisis del Dilema Ético Profesional}
\subsection{Descripción de la Situación Dilemática}
Un recién egresado contratado como asesor técnico en una empresa comercializadora de hortalizas en el Valle del Pilón se enfrenta a una infestación severa de trips y mosquita blanca previa a la cosecha. La gerencia técnica le ordena prescribir un insecticida organofosforado de alta residualidad no autorizado para esa hortaliza, alegando que de no aplicarse se perdería la certificación del lote de exportación y su permanencia laboral.
\subsection{Conflictos de Valores Identificados}
\begin{enumerate}
    \item \textbf{Seguridad alimentaria y salud pública vs. Beneficio económico privado:} La presencia de residuos químicos tóxicos en alimentos atenta contra la vida de los consumidores finales.
    \item \textbf{Legalidad fitosanitaria vs. Convivencia corporativa:} Cumplir la Ley Federal de Sanidad Vegetal frente a la presión jerárquica empresarial.
    \item \textbf{Autonomía deontológica vs. Vulnerabilidad laboral:} Mantener el criterio técnico científico frente a amenazas de despido.
\end{enumerate}

\section{Marco Deontológico e Institucional UANL}
El \textcite{uanl_codigo2021} establece con claridad que los egresados de la Universidad deben regirse por los principios rectores de:
\begin{itemize}
    \item \textbf{Integridad e Integridad Científica:} Actuar con verdad, honestidad y congruencia, subordinando el interés particular al bienestar social.
    \item \textbf{Responsabilidad Social y Ecológica:} Promover un desarrollo ambientalmente viable, rechazando prácticas agronómicas ecocidas o nocivas para la salud de jornaleros y consumidores.
    \item \textbf{Cultura de la Legalidad:} Defender las leyes e instituciones de la República, absteniéndose de actos de simulación, soborno o falsedad en bitácoras de campo \parencite{cndh2020}.
\end{itemize}

\section{Matriz de Práctica Reflexiva y Toma de Decisiones}
Siguiendo el modelo deliberativo de \textcite{giusti2018}, se evalúan tres cursos de acción posibles:

\begin{table}[htbp]
\centering
\small
\caption{Evaluación reflexiva de cursos de acción ante el dilema fitosanitario.}
\label{tab:dilema-agroquimicos}
\begin{tabularx}{\textwidth}{@{}l Y Y Y@{}}
\toprule
\textbf{Curso de Acción} & \textbf{Consecuencias Inmediatas} & \textbf{Impacto Ético-Legal} & \textbf{Veredicto} \\
\midrule
\textbf{Opción A:} Ceder a la orden empresarial. & Conservación del empleo temporal; control de plaga. & Fraude fitosanitario, delito contra la salud, complicidad en ilícito. & \textbf{Rechazada:} Quebranta la legalidad y el código deontológico. \\
\addlinespace
\textbf{Opción B:} Renuncia silenciosa sin intervenir. & Evita la implicación directa personal. & Omisión culposa; abandono moral de los consumidores que consumirán producto tóxico. & \textbf{Rechazada:} Falta de responsabilidad cívica y valor civil. \\
\addlinespace
\textbf{Opción C:} Propuesta técnica fundada de biocontrol y negativa escrita. & Presenta alternativas de bajo impacto (extractos botánicos, jabón potásico y hongos entomopatógenos) y sustenta la prohibición legal. & Alineada con el Estado de derecho, la ética del cuidado y la dignidad profesional. & \textbf{Adoptada:} Acción ética legítima y transformadora. \\
\bottomrule
\end{tabularx}
\end{table}

\section{Conclusiones}
1. El verdadero profesional del agro no supedita su ética técnica a presiones corporativas transitorias; la protección de la salud humana y el respeto a la legalidad fitosanitaria son bienes no negociables.
2. La formación ética universitaria dota al agrónomo de las herramientas reflexivas y argumentativas para liderar transiciones agroecológicas sustentables bajo la primacía del Estado de derecho.

\newpage
\printbibliography[title={Referencias bibliográficas}]
\end{document}

